% -----------------------------------------------------------------------------
% Fourier, distribution and contour-analysis foundations before propagators
% -----------------------------------------------------------------------------

传播子的定义反复依赖 Fourier 变换、分布与复分析中的边界值，因此这些数学约定必须先由定义固定。这些工具按依赖关系自然出现：测试函数先给出 Dirac $\delta$ 的定义，主值与 $\ii0^+$ 边界值随后规定奇点怎样取极限，复能量积分的闭合方向最后把这些边界值转化为不同传播条件。后续“推迟”“超前”或“Feynman”传播子只是在这些数学对象上加入不同物理边界条件。

Dirac $\delta$ 不是普通函数，而是作用在测试函数上的线性泛函。对连续测试函数 $f$，其平移形式的定义是
\begin{equation}
 \boxed{
 \int\dd y\,\delta(x-y)f(y)=f(x).}
 \label{eq:delta-def-action-r13}
\end{equation}
与全书 Fourier 号规相容的常用表示为
\begin{equation}
 \boxed{
 \delta(t-t')
 =\int\frac{\dd\omega}{2\pi}\ee^{-\ii\omega(t-t')},
 \qquad
 \delta^{(3)}(\mathbf r-\mathbf r')
 =\int\frac{\dd^3 p}{(2\pi\hbar)^3}
 \ee^{+\ii\mathbf p\cdot(\mathbf r-\mathbf r')/\hbar}.}
 \label{eq:delta-fourier-representation-dp68}
\end{equation}
质量壳 $\delta$ 函数和长时间 $\delta$ 极限都应通过测试函数作用或变量替换理解，而不能把 $\delta$ 当成在一点取无穷大的普通函数。

Heaviside 分布的导数由分布导数定义固定：
\begin{equation}
 \boxed{\frac{\dd\Theta(x)}{\dd x}=\delta(x).}
 \label{eq:heaviside-delta-dp68}
\end{equation}
因此时间排序对象求导时会产生等时接触项。若 $F(t)$ 在 $t=t'$ 连续可微，则
\begin{equation}
 \boxed{
 \partial_t[\Theta(t-t')F(t)]
 =\delta(t-t')F(t')
 +\Theta(t-t')\partial_tF(t).}
 \label{eq:theta-product-derivative-dp68}
\end{equation}

复平面边界值从有限 $\epsilon>0$ 开始。对实变量 $x$，
\begin{equation}
 \boxed{
 \frac1{x\pm\ii\epsilon}
 =\frac{x}{x^2+\epsilon^2}
 \mp\ii\frac{\epsilon}{x^2+\epsilon^2}.}
 \label{eq:plemelj-split-r13}
\end{equation}
这里 $\epsilon$ 与 $x$ 具有相同量纲。主值分布定义为
\begin{equation}
 \boxed{
 \left\langle\operatorname{PV}\frac1x,f\right\rangle
 \equiv\lim_{\epsilon\to0^+}
 \left(\int_{-\infty}^{-\epsilon}+\int_{\epsilon}^{\infty}\right)
 \frac{f(x)}x\,\dd x.}
 \label{eq:principal-value-definition-dp68}
\end{equation}
在分布意义下取 $\epsilon\to0^+$，正文\eqnrefs{eq:plemelj-split-r13} 变成
\begin{equation}
 \boxed{
 \frac1{x\pm\ii0^+}
 =\operatorname{PV}\frac1x\mp\ii\pi\delta(x).}
 \label{eq:sokhotski-plemelj-dp68}
\end{equation}
这条边界值恒等式直接控制传播子的虚部与在壳 $\delta$ 项。

轮廓闭合方向由 Fourier 指数在大半圆上的衰减决定。令 $z=\operatorname{Re}z+\ii\operatorname{Im}z$，取实时间差 $\tau$，则
\begin{equation}
 \boxed{
 \left|\ee^{-\ii z\tau/\hbar}\right|
 =\exp\!\left(\frac{\tau\,\operatorname{Im}z}{\hbar}\right).}
 \label{eq:jordan-exponential-r13}
\end{equation}
因此 $\tau>0$ 时向下闭合、$\tau<0$ 时向上闭合。若
$I(\tau)=\int_{-\infty}^{\infty}\dd z\,F(z)\ee^{-\ii z\tau/\hbar}/(2\pi)$
且半圆积分消失，则取向进一步给出
\begin{equation}
 \boxed{
 \begin{aligned}
 \tau>0:&\quad
 I(\tau)=-\ii\sum_{z_j\in\mathbb C_-}
 \operatorname{Res}_{z_j}[F(z)\ee^{-\ii z\tau/\hbar}],\\
 \tau<0:&\quad
 I(\tau)=+\ii\sum_{z_j\in\mathbb C_+}
 \operatorname{Res}_{z_j}[F(z)\ee^{-\ii z\tau/\hbar}].
 \end{aligned}}
 \label{eq:contour-orientation-residue-dp68}
\end{equation}
下半圆是顺时针，因此第一行多一个负号；这与“选哪些极点”是两个不同的问题。

若 $z_0$ 是 $m$ 阶极点，可局部写成
\begin{equation}
 \boxed{F(z)=\frac{g(z)}{(z-z_0)^m},\qquad g(z)\ \text{在 }z_0\text{ 解析}.}
 \label{eq:higher-pole-local-form-dp68}
\end{equation}
其留数为
\begin{equation}
 \boxed{
 \operatorname{Res}_{z=z_0}F(z)
 =\frac1{(m-1)!}
 \left.\frac{\dd^{m-1}}{\dd z^{m-1}}
 [(z-z_0)^mF(z)]\right|_{z=z_0}.}
 \label{eq:higher-pole-residue-dp68}
\end{equation}
